\documentclass[11pt]{article}
\usepackage[margin=1in]{geometry}
\usepackage{amsmath,amssymb,mathtools}
\usepackage{booktabs}
\usepackage{enumitem}
\usepackage{hyperref}
\usepackage{xcolor}
\usepackage{algorithm}
\usepackage{algpseudocode}

\title{Density-Ratio VGB Decoding ($r=1$)}
\author{}
\date{}

\newcommand{\pb}{\pi_{\mathrm{big}}}
\newcommand{\ps}{\pi_{\mathrm{small}}}
\newcommand{\Vocab}{\mathcal{V}}
\newcommand{\Y}{\mathcal{Y}}
\newcommand{\Z}{\mathcal{Z}}
\newcommand{\R}{R}
\newcommand{\mask}{\mathtt{[mask]}}
\newcommand{\Wt}{\widetilde W}

\begin{document}
\maketitle

\noindent\textbf{Reference paper:}
\href{https://arxiv.org/pdf/2606.28301}{VGB for Masked Diffusion Model: Efficient Test-time Scaling for Reward Satisfaction and Sample Editing}.

\paragraph{Scope.}
This note specializes Definition 4.2 of the paper to singleton moves ($r=1$) and to a density-ratio construction in which a small model is the reference model and a larger model is the target/verifier model. We assume a fixed prompt $x$ is given before decoding. The prompt is never masked or edited; all states below refer only to the $n$ generated positions following the prompt.

\section{Setup and notation}

Let $\Vocab$ be the vocabulary and let
\begin{equation}
\Y = \Vocab^n,
\qquad
\Z = (\Vocab\cup\{\mask\})^n.
\end{equation}
A state $z\in\Z$ is the current generated suffix, with some positions revealed and the rest masked. Define
\begin{equation}
\R(z) \coloneqq \{i\in[n]: z_i\neq\mask\},
\qquad
k \coloneqq |\R(z)|.
\end{equation}
Thus $k$ is the current depth of the walk, $n-k$ positions remain masked, and $z\in\Y$ exactly when $k=n$.

For any masked position $i\notin\R(z)$ and token $a\in\Vocab$, define the forward child
\begin{equation}
c_{i,a} \coloneqq z^{i\leftarrow a}.
\end{equation}
For any revealed position $i\in\R(z)$, define the backward parent
\begin{equation}
p_i \coloneqq z^{-i},
\end{equation}
obtained by replacing $z_i$ with $\mask$.

We use two prompt-conditioned sequence models:
\begin{equation}
\ps(\cdot\mid x)
\quad\text{and}\quad
\pb(\cdot\mid x),
\end{equation}
where the small model is the VGB reference model and the big model is the desired target model.

For $m\in\{\mathrm{small},\mathrm{big}\}$, the probability of a partial state is the compatible-completion marginal
\begin{equation}
\pi_m(z\mid x)
\coloneqq
\sum_{y\in\mathcal C(z)} \pi_m(y\mid x),
\end{equation}
where
\begin{equation}
\mathcal C(z)
\coloneqq
\{y\in\Y : y_j=z_j \text{ for all }j\in\R(z)\}.
\end{equation}
For a child $c_{i,a}$, the corresponding one-token conditional is
\begin{equation}
\pi_m(Y_i=a\mid x,z)
=
\frac{\pi_m(c_{i,a}\mid x)}{\pi_m(z\mid x)}.
\label{eq:conditional}
\end{equation}

\section{Density-ratio value}

Choose the terminal reward
\begin{equation}
\tau(x,y)
\coloneqq
\frac{\pb(y\mid x)}{\ps(y\mid x)}.
\end{equation}
Then the VGB reward-tilted target is exactly the big-model distribution:
\begin{equation}
\pi^\star(y\mid x)
\propto
\ps(y\mid x)\tau(x,y)
=
\pb(y\mid x).
\end{equation}
Since $\pb(\cdot\mid x)$ is normalized, the normalizing constant is one.

The corresponding ideal VGB value of a partial state is
\begin{align}
V(z)
&\coloneqq
\mathbb E_{Y\sim\ps(\cdot\mid x)}
\left[\tau(x,Y)\mid Y\in\mathcal C(z)\right] \\
&=
\frac{\pb(z\mid x)}{\ps(z\mid x)}.
\label{eq:value-density-ratio}
\end{align}
The absolute partial-state marginals in \eqref{eq:value-density-ratio} are generally intractable. The key observation is that Definition 4.2 only requires them through ratios that can be expressed using local model conditionals.

\subsection{Neighboring-state ratios}

For a forward child $c_{i,a}$, equation \eqref{eq:conditional} gives
\begin{equation}
\frac{V(c_{i,a})}{V(z)}
=
\frac{\pb(Y_i=a\mid x,z)}{\ps(Y_i=a\mid x,z)}.
\label{eq:forward-value-ratio}
\end{equation}
Equivalently,
\begin{equation}
V(c_{i,a})
=
V(z)
\frac{\pb(Y_i=a\mid x,z)}{\ps(Y_i=a\mid x,z)}.
\end{equation}

For a backward parent $p_i=z^{-i}$, the same identity applied to the edge $p_i\to z$ gives
\begin{equation}
\frac{V(p_i)}{V(z)}
=
\frac{\ps(Y_i=z_i\mid x,p_i)}{\pb(Y_i=z_i\mid x,p_i)}.
\end{equation}
Define this tractable leave-one-out ratio as
\begin{equation}
\boxed{
\rho_i(z)
\coloneqq
\frac{V(p_i)}{V(z)}
=
\frac{\ps(Y_i=z_i\mid x,z^{-i})}
     {\pb(Y_i=z_i\mid x,z^{-i})}
}.
\label{eq:rho}
\end{equation}
Large $\rho_i$ means the small model favors the committed token $z_i$ more strongly than the big model does in the current leave-one-out context.

\section{Definition 4.2 specialized to $r=1$}

For singleton moves, the depth coefficient is
\begin{equation}
\boxed{
 s_{k,1} = \binom{n-1}{k}^{-1}
}
\end{equation}
and satisfies the balancing identity
\begin{equation}
(n-k)s_{k,1}=k s_{k-1,1}.
\end{equation}

Consider an interior state $z\notin\Y$, so $0\le k<n$. For a non-leaf forward child $c_{i,a}\notin\Y$, Definition 4.2 gives
\begin{equation}
w_{\mathrm{fwd}}(z\to c_{i,a})
=
s_{k,1}
V(z)^\lambda
V(c_{i,a})
\ps(Y_i=a\mid x,z).
\label{eq:def42-forward}
\end{equation}
For a backward parent $p_i=z^{-i}$,
\begin{equation}
w_{\mathrm{bwd}}(z\to p_i)
=
s_{k-1,1}
V(z)
V(p_i)^\lambda.
\label{eq:def42-backward}
\end{equation}

The paper normalizes all outgoing edge weights by
\begin{equation}
W(z)=F(z)+B(z),
\end{equation}
where
\begin{equation}
F(z)=\sum_{i\notin\R(z)}\sum_{a\in\Vocab}
w_{\mathrm{fwd}}(z\to c_{i,a}),
\qquad
B(z)=\sum_{i\in\R(z)}w_{\mathrm{bwd}}(z\to p_i).
\end{equation}

\section{Reduced transition weights}

\subsection{Forward edge}

Substitute \eqref{eq:forward-value-ratio} into \eqref{eq:def42-forward}:
\begin{align}
w_{\mathrm{fwd}}(z\to c_{i,a})
&=
s_{k,1}V(z)^\lambda
\left[
V(z)\frac{\pb(a\mid x,z)}{\ps(a\mid x,z)}
\right]
\ps(a\mid x,z) \\
&=
\boxed{
s_{k,1}V(z)^{\lambda+1}\pb(Y_i=a\mid x,z)
}.
\label{eq:reduced-forward-edge}
\end{align}
Thus the reference-model token probability cancels locally.

For a \emph{fixed} masked coordinate $i$, sum over all token values:
\begin{align}
\sum_{a\in\Vocab} w_{\mathrm{fwd}}(z\to c_{i,a})
&=
s_{k,1}V(z)^{\lambda+1}
\sum_{a\in\Vocab}\pb(Y_i=a\mid x,z) \\
&=
s_{k,1}V(z)^{\lambda+1},
\end{align}
because $\pb(Y_i=\cdot\mid x,z)$ is a normalized categorical distribution. This holds separately for every masked coordinate $i$; there is no random choice of $i$ in this summation.

There are exactly $n-k$ masked coordinates, so
\begin{equation}
\boxed{
F(z)
=(n-k)s_{k,1}V(z)^{\lambda+1}
}.
\label{eq:forward-mass}
\end{equation}

\subsection{Backward edge}

Using $V(p_i)=V(z)\rho_i$, equation \eqref{eq:def42-backward} becomes
\begin{align}
w_{\mathrm{bwd}}(z\to p_i)
&=
s_{k-1,1}V(z)\left[V(z)\rho_i\right]^\lambda \\
&=
\boxed{
s_{k-1,1}V(z)^{\lambda+1}\rho_i^\lambda
}.
\label{eq:reduced-backward-edge}
\end{align}
Hence
\begin{equation}
\boxed{
B(z)
=
s_{k-1,1}V(z)^{\lambda+1}
\sum_{i\in\R(z)}\rho_i^\lambda
}.
\label{eq:backward-mass}
\end{equation}

\subsection{Reduced normalizer}

Combining \eqref{eq:forward-mass} and \eqref{eq:backward-mass},
\begin{equation}
W(z)
=
V(z)^{\lambda+1}
\left[
(n-k)s_{k,1}
+
s_{k-1,1}\sum_{i\in\R(z)}\rho_i^\lambda
\right].
\end{equation}
The absolute, intractable factor $V(z)^{\lambda+1}$ is common to every outgoing transition and therefore cancels during normalization. Define
\begin{equation}
\boxed{
\Wt(z)
\coloneqq
(n-k)s_{k,1}
+
s_{k-1,1}\sum_{i\in\R(z)}\rho_i^\lambda
}.
\label{eq:reduced-W}
\end{equation}
No running estimate or recursion for the absolute value $V(z)$ is needed by the sampler.

The direction probabilities are therefore
\begin{align}
\boxed{
P(\mathrm{forward}\mid z)
=
\frac{(n-k)s_{k,1}}{\Wt(z)}
},
\label{eq:p-forward}\\[4pt]
\boxed{
P(\mathrm{backward}\mid z)
=
\frac{s_{k-1,1}\sum_{i\in\R(z)}\rho_i^\lambda}{\Wt(z)}
}.
\label{eq:p-backward}
\end{align}

Conditional on a forward move, equation \eqref{eq:reduced-forward-edge} gives
\begin{equation}
\boxed{
P(i,a\mid \mathrm{forward},z)
=
\frac{1}{n-k}\pb(Y_i=a\mid x,z)
},
\end{equation}
so the implementation can sample
\begin{equation}
i\sim\mathrm{Unif}([n]\setminus\R(z)),
\qquad
a\sim\pb(Y_i=\cdot\mid x,z).
\label{eq:forward-sampling}
\end{equation}

Conditional on a backward move, equation \eqref{eq:reduced-backward-edge} gives
\begin{equation}
\boxed{
P(i\mid \mathrm{backward},z)
=
\frac{\rho_i^\lambda}
{\sum_{j\in\R(z)}\rho_j^\lambda}
},
\qquad i\in\R(z).
\label{eq:backward-sampling}
\end{equation}
The selected token is then remasked.

\section{Model evaluations per walk step}

For the exact full-neighborhood $r=1$ rule, construct every leave-one-out parent
\begin{equation}
\mathcal P(z)=\{z^{-i}:i\in\R(z)\}.
\end{equation}
For each row $p_i=z^{-i}$, we need only two scalar probabilities:
\begin{equation}
\ps(Y_i=z_i\mid x,p_i),
\qquad
\pb(Y_i=z_i\mid x,p_i),
\end{equation}
which determine $\rho_i$ through \eqref{eq:rho}.

These $k$ parent states can be evaluated as a batch. A convenient implementation performs:
\begin{enumerate}[leftmargin=2em]
\item one batched call to the small model on the parent batch;
\item one batched call to the big model on the same parent batch, optionally including the current state $z$ as an additional row so that the big-model forward distribution is already available if a forward move is selected.
\end{enumerate}
Thus the algorithm can use two \emph{batched} model calls per Markov step, but the batch contains $k$ (or $k+1$) sequences. This is not equivalent in compute cost to two single-sequence forward passes.

If $k=0$, there are no backward parents and the move is necessarily forward. If $k=n$, the state is fully revealed. The pseudocode below uses first-leaf stopping and returns the completion; the exact VGB Markov chain may instead continue by applying the leaf branch of Definition 4.2.

\section{Decoding algorithm}

\begin{algorithm}[H]
\caption{Density-Ratio VGB Decoding, Singleton Moves ($r=1$)}
\label{alg:density-ratio-vgb}
\begin{algorithmic}[1]
\Require prompt $x$; small model $\ps$; big model $\pb$; generation length $n$; remasking parameter $\lambda\ge 0$
\State $z\gets(\mask,\ldots,\mask)\in\Z$
\While{$|\R(z)|<n$}
    \State $k\gets|\R(z)|$
    \If{$k=0$}
        \State $d\gets\mathrm{forward}$
    \Else
        \ForAll{$i\in\R(z)$}
            \State $p_i\gets z^{-i}$
        \EndFor
        \State Evaluate all $p_i$ as a batch under $\ps$ and $\pb$
        \ForAll{$i\in\R(z)$}
            \State $\rho_i\gets \ps(Y_i=z_i\mid x,p_i)\,/\,\pb(Y_i=z_i\mid x,p_i)$
        \EndFor
        \State $s_f\gets\binom{n-1}{k}^{-1}$
        \State $s_b\gets\binom{n-1}{k-1}^{-1}$
        \State $M_f\gets(n-k)s_f$
        \State $M_b\gets s_b\sum_{i\in\R(z)}\rho_i^\lambda$
        \State $\Wt\gets M_f+M_b$
        \State Sample $d\in\{\mathrm{forward},\mathrm{backward}\}$ with probabilities $M_f/\Wt$ and $M_b/\Wt$
    \EndIf

    \If{$d=\mathrm{forward}$}
        \State Sample $i\sim\mathrm{Unif}([n]\setminus\R(z))$
        \State Evaluate $\pb(Y_i=\cdot\mid x,z)$ if it was not already computed
        \State Sample $a\sim\pb(Y_i=\cdot\mid x,z)$
        \State $z\gets z^{i\leftarrow a}$
    \Else
        \State Sample $i\in\R(z)$ with probability $\rho_i^\lambda/\sum_{j\in\R(z)}\rho_j^\lambda$
        \State $z\gets z^{-i}$
    \EndIf
\EndWhile
\State \Return the prompt-conditioned completion $(x,z)$
\end{algorithmic}
\end{algorithm}

\section{Implementation interpretation}

At an interior state, the complete transition rule can be summarized as
\begin{equation}
\boxed{
\begin{array}{ll}
\text{Forward mass:} & (n-k)s_{k,1},\\[3pt]
\text{Backward mass at position }i: & s_{k-1,1}\rho_i^\lambda,\\[3pt]
\rho_i: & \displaystyle\frac{\ps(z_i\mid x,z^{-i})}{\pb(z_i\mid x,z^{-i})}.
\end{array}}
\end{equation}
The common factor $V(z)^{\lambda+1}$ from Definition 4.2 never needs to be evaluated. Forward moves choose a masked coordinate uniformly and draw its token from the big model; backward moves preferentially erase revealed tokens whose leave-one-out probability is relatively larger under the small model than under the big model.

\paragraph{Model-consistency caveat.}
The density-ratio derivation is exact when the masked conditionals of each model are compatible with an underlying coherent sequence-level distribution, as assumed in the paper's probabilistic formulation. For practical masked language models this may hold only approximately; in that case the local ratios in \eqref{eq:rho} define an approximate density-ratio VGB implementation.

\end{document}
